\documentclass[10pt,a4paper]{article}
\usepackage{amsmath,amssymb,amsthm,graphicx}
% Minimal retained preamble for S3 SIGMA 2022 article 023 cards taken
% from sources/sigma-2022-023-post-lie-magnus-expansion-and-bch-recursion/
% pinned-original-source.tex (doi:10.3842/SIGMA.2022.023, SIGMA 18 (2022),
% CC BY 4.0).
%
% Omitted and why: the publisher documentclass{sigma} block (source
% lines 1-7 with its ifpdf branch) - the sigma class is unavailable in
% the toolchain and the assembler supplies the article class
% (unavailable-class precedent). The \usepackage{wasysym} line (source
% line 9): no delivered span of any card of this batch uses a wasysym-only
% symbol (the \rhd of the delivered spans is the amssymb relation, and
% the author's wasysym-dependent \rightarrowRHD/\longrightarrowRHD
% arrows, source lines 13-14, are unused by every delivered span).
% \lbutcher (source line 12) is unused by every delivered span. Of the
% author macro block only the declarations used by the delivered spans
% are copied verbatim below, in source order: the \nc shorthand (source
% line 11, needed because the copied macro lines are written with it),
% \Cal (line 41), \mop (line 42), \wt (line 46; the delivered spans use
% the LaTeX-standard \widehat/\widetilde directly, so the author's
% \wh shorthand at line 47 is not needed). Omitted as unused by every
% delivered span: \vep/\ep (29), \sigmat (30), \ostar (31), \grad (33),
% \fil (34), the blackboard-letter block \BA-\ZZ (36-40), \mopl (43),
% \frakg (44), \g (45), \wh (47), \un (48), \fleche (50) and any
% further macro-block lines. The author equation numbering (source line
% 16) and the author theorem declarations (source lines 18-27, master
% counter Theorem sectionwise with Corollary/Lemma/Proposition hanging on
% it, the starred Theorem* at line 19, and the definition-style group
% with its brace pair on lines 23 and 27) are copied verbatim below. No
% per-card \usepackage additions are needed (CONTROL 8.8.162: nothing
% outside the amsmath/amssymb/amsthm baseline occurs in the delivered
% spans).
\usepackage{amsmath,amssymb,amsthm}
\numberwithin{equation}{section}

\newtheorem{Theorem}{Theorem}[section]
\newtheorem*{Theorem*}{Theorem}
\newtheorem{Corollary}[Theorem]{Corollary}
\newtheorem{Lemma}[Theorem]{Lemma}
\newtheorem{Proposition}[Theorem]{Proposition}
 { \theoremstyle{definition}
\newtheorem{Definition}[Theorem]{Definition}
\newtheorem{Note}[Theorem]{Note}
\newtheorem{Example}[Theorem]{Example}
\newtheorem{Remark}[Theorem]{Remark} }

\newcommand{\nc}{\newcommand}
\nc{\Cal}[1]{{\mathcal {#1}}}
\nc{\mop}[1]{\mathop{\hbox {\rm #1} }}
\nc{\wt}{\widetilde}
\title{Post-Lie Magnus expansion equals the weighted BCH-recursion at weight one (Theorem 5.3, the paper's main theorem)}
\author{Al-Kaabi, Mahdi J. Hasan, Ebrahimi-Fard, Kurusch, Manchon, Dominique}
\date{doi:10.3842/SIGMA.2022.023; CC BY 4.0}
\begin{document}
\maketitle
\setcounter{section}{5}
\setcounter{Theorem}{2}
\setcounter{equation}{0}
\noindent\textit{Editorial note.} All delivered mathematical text is the author's own, copied from the pinned source file sources/sigma-2022-023-post-lie-magnus-expansion-and-bch-recursion/pinned-original-source.tex; no statement or proof step was rewritten, reordered or removed.\par\medskip
\noindent\textit{Editorial note.} Printed numbers anchored to the published PDF (SIGMA 18 (2022) 023, stream-extracted text lane-s3/sig22-pdf-anchor-023-streams.txt, census pdf\_anchor\_evidence lane-s3/sig22-pdf-anchor-023-summary.json): Theorem 5.1 (T9) and Theorem 5.3 ('main Theorem') surface in the stream text ('Proof of Theorem 5.1', s021/s022); Proposition 5.2 is bracketed by the confirmed neighbours; the sectionwise shared master counter Theorem is declared in the author's newtheorem block (source lines 18-27: Corollary, Lemma, Proposition hang on the Theorem counter, the starred Theorem* at line 19, the definition-style group with its brace pair on lines 23 and 27).\par\medskip
\noindent\textit{Editorial note.} Theorem-family numbers and equation numbers inside the card are locally renumbered (contexts taken from other sections of the source cannot keep their paper numbers under the sectionwise counters); every cross-reference is a symbolic ref and resolves inside the card; cite closure delivered as card-local bibliography fragments with href rendered plain if any (no hyperref in baseline); counters seeded from in-section source counts via the master counter name; author macro subset copied verbatim per CONTROL 8.8.165 (lines 11/41/42/46); no extra packages needed (wasysym omitted: no wasysym-only symbol in any delivered span).\par\medskip
\noindent\textit{Editorial note.} Key-shape watch: the claim key of the main-theorem card carries the verbatim author label 'main Theorem' including its space (source line 674); this is deliberate and stream-anchored.\par\medskip
\noindent\textit{Editorial note.} The counted span is the author Theorem environment with the verbatim spaced label 'main Theorem' (source lines 674-676) with its single proof environment (source lines 678-694).\par\medskip
\noindent\textit{Editorial note.} The statement's \textbackslash{}eqref\{pl-Magnus\} and \textbackslash{}eqref\{chi-lambda-rec\} and the proof's \textbackslash{}eqref\{exp-moins-lambda-x\}, \textbackslash{}eqref\{IPLME\} and \textbackslash{}eqref\{ET9\} all resolve through the delivered contexts; the context-span \textbackslash{}cite keys (KLM06 at source line 416, ELM2015 at source line 522) are closed by the card-local bibliography fragment.\par\medskip
\noindent\textit{Editorial note.} Printed numbers anchored to the published PDF (SIGMA 18 (2022) 023, stream-extracted text lane-s3/sig22-pdf-anchor-023-streams.txt, census pdf\_anchor\_evidence lane-s3/sig22-pdf-anchor-023-summary.json): Theorem 5.3 confirmed.\par\medskip
\noindent\textit{Editorial scope.} Single environment card cut from the SIGMA article doi:10.3842/SIGMA.2022.023 (sigma-2022-023-post-lie-magnus-expansion-and-bch-recursion); statement and proof are the author's own text at source lines 674-676 and 678-694.\par\medskip
\noindent\textit{Allowed standard prerequisites (editorial).}\begin{itemize}
\item Weight-one specialisation of the inverse weighted recursion. Delivered in context (source lines 416-440): exp-moins-lambda-x with its KLM06 quotation frame.. Load-bearing: the proof opens from \textbackslash{}eqref\{exp-moins-lambda-x\} specialised to lambda=1 (source line 679).
\item Inverse post-Lie Magnus expansion IPLME. Delivered in context (source lines 554-559).. Load-bearing: combined with ET9 and res.2 at source line 689.
\item Weight-one exponential factorization ET9 (Theorem 5.1). Delivered in context (source lines 605-611).. Load-bearing: \textbackslash{}eqref\{ET9\} enters the res.3 combination (source line 689).
\end{itemize}\medskip
\begin{Theorem}\label{main Theorem}
The post-Lie Magnus expansion $\chi$, described in \eqref{pl-Magnus}, coincides with the weighted BCH-recursion $\chi_{\lambda}$ recursively given by \eqref{chi-lambda-rec}, with weight $\lambda=1$.
\end{Theorem}

\begin{Proposition}[{\cite[Proposition 11]{KLM06}}]
Let $\big(\Cal{A}, \Cal{R}\big)$ be a complete filtered Rota--Baxter algebra of weight $\lambda\neq 0$, and set $\widetilde{\Cal{R}} := -\lambda {\rm id}_{\Cal{A}}- \Cal{R}$. The $\lambda$-weighted BCH-recursion is written
\begin{equation}
\label{chi-lambda-rec}
\chi_{\lambda}(x)
= x + \frac{1}{\lambda} \wt{{\rm BCH}}\big(\Cal{R}(\chi_{\lambda}(x)),
\widetilde{\Cal{R}}(\chi_{\lambda}(x))\big),
\end{equation}
for all $x\in\Cal A_1$. It can be simplified to
\begin{equation*}
\chi_{\lambda}(x)
= x - \frac{1}{\lambda} \wt{{\rm BCH}}\big({-}\Cal{R}(\chi_{\lambda}(x)), \lambda x\big).
\end{equation*}
Its inverse is given by
\begin{equation*}
\chi_{\lambda}^{-1}(x)
= x - \frac{1}{\lambda} \wt{{\rm BCH}}\big(\Cal{R}(x), \widetilde{\Cal{R}}(x)\big).
\end{equation*}
Moreover, the factorization obtained in \eqref{exp-fact} becomes
\begin{equation}
\label{exp-moins-lambda-x}
\exp\big(\Cal{R}(\chi_{\lambda}(x))\big)\exp\big(\widetilde{\Cal{R}}
(\chi_{\lambda}(x))\big) = \exp(-\lambda x).
\end{equation}
\end{Proposition}

\begin{equation}\label{exp-fact}
\exp\big(\Cal{R}(\chi(x))\big)\exp\big(\widetilde{\Cal{R}}(\chi(x))\big) = \exp(x),
\end{equation}

Both these exponential functions map $\mop{Prim}\big(F^{\rm pl}_E\big)$ bijectively onto $G\big(F^{\rm pl}_E\big)$. See \cite{ELM2015} for details. The post-Lie Magnus expansion $\chi$ is the bijective map from $\widehat{\Cal{L}\big(\Cal{T}^{\rm pl}_E\big)}$ onto itself defined by
\begin{displaymath}
\exp^{*}\big(\chi(f)\big) = \exp(f),
\end{displaymath}
namely,
\begin{equation}
\label{pl-Magnus}
\chi(f) = \log^{\ast}\big(\exp(f)\big).
\end{equation}

\subsection{Inverse post-Lie Magnus expansion}

The inverse post-Lie Magnus expansion $\theta$ is the bijective map from $\widehat{\Cal{L}\big(\Cal{T}^{\rm pl}_E\big)}$ onto itself given by the following formula
\begin{equation}\label{IPLME}
\theta(f) = \log(\exp^*(f))
\end{equation}

\begin{Theorem}\label{T9}
Let $\big(\Cal{A}, \Cal{R}\big)$ be a complete filtered Rota--Baxter algebra of weight $\lambda = 1$. We have the following equality in $\widehat{\Cal U(\Cal A)}$ for any $x\in\Cal A$ and $t \in K$:
\begin{equation}\label{ET9}
\exp^{\ast}(tx) = \exp\big({-}t \widetilde{\Cal{R}}(x)\big) \exp\big({-}t \Cal{R}(x)\big),
\end{equation}
where $\widetilde{\Cal{R}}= -{\rm id}_{\Cal{A}}-\Cal{R}$, and $\ast$ is the associative product defined in \eqref{ast}, using the post-Lie product $x \rhd y = [\Cal{R}(x), y]$.
\end{Theorem}

\begin{equation}\label{ast}
X \ast Y := X_{(1)}\big(X_{(2)} \rhd Y\big),
\end{equation}

\begin{proof}
From equation \eqref{exp-moins-lambda-x}, specialized to $\lambda=1$, and by setting $\theta_{\rm BCH} :=\chi_{1}^{-1}$, we obtain that
\begin{displaymath}
\exp\big({-}\theta_{\rm BCH}(tx)\big)
= \exp\big(t \Cal{R}(x)\big) \exp\big(t \widetilde{\Cal{R}}(x)\big)
\end{displaymath}
in $\widehat{\Cal U(\Cal A)}$, which is equivalent to
\begin{equation}\label{res.2}
\exp\big(\theta_{\rm BCH}(tx)\big)
= \exp\big({-}t \widetilde{\Cal{R}}(x)\big) \exp\big({-}t \Cal{R}(x)\big).
\end{equation}
From \eqref{IPLME}, \eqref{ET9} and \eqref{res.2} we have
\begin{equation}\label{res.3}
\exp\big(\theta_{\rm BCH}(tx)\big) = \exp^{\ast}(tx)= \exp\big(\theta(tx)\big).
\end{equation}
Then the two $\theta$'s, namely the inverse BCH-recursion in~\eqref{res.2} and the inverse post-Lie Magnus expansion \eqref{res.3}, do coincide.
\end{proof}

\begin{thebibliography}{99}
% Card-local bibliography fragment: the \cite keys inside the
% delivered spans of this card (here: context-span cites).
% Adapted from the author's thebibliography block (source lines
% 761-914) of sources/sigma-2022-023-post-lie-magnus-expansion-
% and-bch-recursion/pinned-original-source.tex, with the \href
% wrappers rendered as plain text (inner content kept) if any, because the card
% baseline carries no hyperref (bibliography typesetting only;
% no statement or proof text is touched).

\bibitem{KLM06}
Ebrahimi-Fard K., Guo L., Manchon D., Birkhoff type decompositions and the
 {B}aker--{C}ampbell--{H}ausdorff recursion, \textit{Comm. Math. Phys.}
 \textbf{267} (2006), 821--845, arXiv:math-ph/0602004.


\bibitem{ELM2015}
Ebrahimi-Fard K., Lundervold A., Munthe-Kaas H.Z., On the {L}ie enveloping
 algebra of a post-{L}ie algebra, \textit{J.~Lie Theory} \textbf{25} (2015),
 1139--1165, arXiv:1410.6350.

\end{thebibliography}
\end{document}
